\chapter{Filas e threads: producer/consumer}

\section{Por que usar fila?}

Parser e banco podem ter velocidades diferentes. O parser pode produzir registros mais rapido do que o Postgres consegue inserir, ou o banco pode ficar esperando I/O enquanto o parser poderia continuar trabalhando.

A arquitetura producer/consumer separa essas responsabilidades:

\begin{verbatim}
producer thread
    le XML e cria HealthRecord
        -> Queue
consumer thread
    pega HealthRecord e insere em batch no Postgres
\end{verbatim}

\section{Queue}

\code{queue.Queue} e uma fila thread-safe da biblioteca padrao.

\begin{lstlisting}[style=devbutterpython]
from queue import Queue

record_queue: Queue[HealthRecord | None] = Queue(maxsize=10_000)
\end{lstlisting}

\section{Explicando cada nomezinho}

\begin{itemize}
  \item \code{Queue}: fila segura para multiplas threads.
  \item \code{maxsize}: limite da fila. Evita crescimento infinito de memoria.
  \item \code{HealthRecord | None}: a fila aceita registros ou \code{None}.
  \item \code{None}: usado como sentinela para sinalizar fim do trabalho.
  \item \code{backpressure}: quando a fila enche, o producer bloqueia. Isso protege a memoria.
\end{itemize}

\section{Producer}

\begin{lstlisting}[style=devbutterpython]
from pathlib import Path
from queue import Queue
import xml.etree.ElementTree as ET

from src.models.health_record import HealthRecord
from src.parser.xml_parser import clean_tag, parse_record_element


def producer(
    xml_file: Path,
    record_queue: Queue[HealthRecord | None],
) -> None:
    record_id = 0
    context = ET.iterparse(xml_file, events=("end",))

    for _, element in context:
        if clean_tag(element.tag) == "Record":
            record_id += 1
            record = parse_record_element(element, record_id)
            record_queue.put(record)

        element.clear()

    record_queue.put(None)
\end{lstlisting}

\section{Consumer}

\begin{lstlisting}[style=devbutterpython]
from queue import Queue

from src.db.repository import HealthRecordRepository
from src.models.health_record import HealthRecord


def consumer(
    record_queue: Queue[HealthRecord | None],
    repository: HealthRecordRepository,
    batch_size: int,
) -> None:
    batch: list[HealthRecord] = []

    while True:
        record = record_queue.get()

        if record is None:
            break

        batch.append(record)

        if len(batch) >= batch_size:
            repository.insert_batch(batch)
            batch.clear()

    if batch:
        repository.insert_batch(batch)
\end{lstlisting}

\section{Rodando com Thread}

\begin{lstlisting}[style=devbutterpython]
from pathlib import Path
from queue import Queue
from threading import Thread

from src.db.repository import HealthRecordRepository
from src.models.health_record import HealthRecord


def run_threaded_pipeline(
    xml_file: Path,
    repository: HealthRecordRepository,
    batch_size: int,
) -> None:
    record_queue: Queue[HealthRecord | None] = Queue(maxsize=10_000)

    producer_thread = Thread(
        target=producer,
        args=(xml_file, record_queue),
        name="xml-producer",
    )

    consumer_thread = Thread(
        target=consumer,
        args=(record_queue, repository, batch_size),
        name="postgres-consumer",
    )

    producer_thread.start()
    consumer_thread.start()

    producer_thread.join()
    consumer_thread.join()
\end{lstlisting}

\section{Quando thread ajuda?}

Threads ajudam quando existe espera de I/O. No projeto, escrever no Postgres e ler arquivos sao I/O-bound. Enquanto uma thread espera o banco responder, outra pode continuar lendo ou preparando dados.

Mas se o gargalo principal for parse CPU-bound, threads podem nao acelerar muito por causa do GIL.

\section{Experimento}

Compare:

\begin{itemize}
  \item pipeline sincrona;
  \item pipeline com producer/consumer e threads;
  \item batch sizes diferentes;
  \item queue maxsize diferente.
\end{itemize}

Pergunta que voce deve responder: \textbf{thread melhorou por causa de I/O ou nao mudou por causa de CPU/GIL?}
